\documentclass[11pt]{article}

\usepackage[margin=1in]{geometry}
\usepackage[T1]{fontenc}
\usepackage{lmodern}
\usepackage{amsmath,amssymb,amsthm,mathtools}

\newtheorem{theorem}{Theorem}
\newtheorem{lemma}[theorem]{Lemma}
\newtheorem{proposition}[theorem]{Proposition}
\newtheorem{corollary}[theorem]{Corollary}
\newtheorem{remark}{Remark}

\title{A Threshold Law for Embedding-Based Graph Retrieval Under a Gaussian Surrogate Model}
\author{Krasnow, Hunt, Wu, Park}
\date{}

\begin{document}
\maketitle

\section{Setup}

Let $G=(V,E)$ be a graph with mean degree $\bar d > 1$. Fix a query node $q \in V$ and an answer node $a \in V$ at shortest-path distance $h$ from $q$. Let $\mathcal{B}_h(q)$ denote the candidate set searched by the retrieval procedure, and write
\[
M = |\mathcal{B}_h(q)|.
\]
In a locally tree-like graph one often has the heuristic scaling
\[
M \approx \bar d^{\,h},
\]
which we use later to express the threshold in hop count.

Let $f:V\to\mathbb{R}^D$ be an embedding map with unit-normalized outputs, and define the similarity score
\[
S_v = \langle f(q), f(v)\rangle.
\]
The retrieval rule is
\[
\hat a = \arg\max_{v\in \mathcal{B}_h(q)} S_v.
\]
We study the probability that $\hat a = a$ under the Gaussian surrogate model below.

\section{Model Assumptions}

We assume the following.

\begin{enumerate}
    \item[\textbf{A1.}] \textbf{(Candidate-set size.)}
    The candidate set has size $M=M(D)\ge 2$.

    \item[\textbf{A2.}] \textbf{(Noise scores.)}
    For the $M-1$ non-answer candidates, the scores are independent and identically distributed as
    \[
    S_{v_1},\dots,S_{v_{M-1}} \stackrel{\mathrm{iid}}{\sim} \mathcal{N}(0,\sigma_D^2).
    \]

    \item[\textbf{A3.}] \textbf{(Answer score.)}
    The answer score is independent of the noise scores and satisfies
    \[
    S_a \sim \mathcal{N}(\delta_D,\sigma_D^2).
    \]

    \item[\textbf{A4.}] \textbf{(High-dimensional normalization.)}
    We take
    \[
    \sigma_D^2 = \frac{1}{D}.
    \]
    Thus the normalized signal level is
    \[
    \Delta_D := \frac{\delta_D}{\sigma_D} = \delta_D \sqrt{D}.
    \]
\end{enumerate}

\begin{remark}
Assumptions A2--A4 define a surrogate model. They are analytically convenient and capture the extreme-value competition between one signal score and many noise scores. They should not be read as an exact model of real graph embeddings, where dependence, heteroskedasticity, and candidate-selection effects may all matter.
\end{remark}

\section{Main Result}

Our first result is an exact formula for the success probability under the surrogate model.

\begin{proposition}[Exact success probability]\label{prop:exact}
Under A2--A4,
\begin{equation}\label{eq:exact_integral}
P(\hat a = a)
=
\int_{-\infty}^{\infty}
\Phi(x+\Delta_D)^{M-1}\phi(x)\,dx,
\end{equation}
where $\phi$ and $\Phi$ denote the standard normal density and distribution function.
\end{proposition}

The next theorem gives the asymptotic threshold law.

\begin{theorem}[Threshold law]\label{thm:threshold}
Assume A2--A4, and let $M=M(D)\to\infty$ as $D\to\infty$. Define
\[
L_D := \sqrt{2\log M(D)}.
\]
Then:
\begin{enumerate}
    \item If
    \[
    \Delta_D - L_D \to +\infty,
    \]
    then
    \[
    P(\hat a = a)\to 1.
    \]

    \item If
    \[
    \Delta_D - L_D \to -\infty,
    \]
    then
    \[
    P(\hat a = a)\to 0.
    \]
\end{enumerate}
Equivalently, retrieval success is governed asymptotically by the comparison between the normalized signal level $\Delta_D$ and the Gaussian noise maximum scale $\sqrt{2\log M(D)}$.
\end{theorem}

When $M$ is parameterized by hop count through $M \approx \bar d^{\,h}$, the theorem yields the corresponding crossover scale in $h$.

\begin{corollary}[Heuristic hop-count threshold]\label{cor:hstar}
Suppose $M \approx \bar d^{\,h}$ with $\bar d>1$, and suppose $\delta_D \equiv \delta > 0$ is constant in $D$. Then the threshold condition
\[
\Delta_D \approx \sqrt{2\log M}
\]
becomes
\[
\delta\sqrt{D} \approx \sqrt{2h\log\bar d},
\]
so the corresponding hop-count scale is
\begin{equation}\label{eq:hstar}
h^*(D) \approx \frac{\delta^2 D}{2\log\bar d}.
\end{equation}
More precisely:
\begin{itemize}
    \item if
    \[
    \delta\sqrt{D} - \sqrt{2h(D)\log\bar d}\to +\infty,
    \]
    then $P(\hat a=a)\to 1$;

    \item if
    \[
    \delta\sqrt{D} - \sqrt{2h(D)\log\bar d}\to -\infty,
    \]
    then $P(\hat a=a)\to 0$.
\end{itemize}
\end{corollary}

\begin{remark}
Corollary~\ref{cor:hstar} is a consequence of the surrogate model together with the size heuristic $M\approx \bar d^{\,h}$. It should be interpreted as a crossover law, not as a claim about all retrieval procedures or all graph/embedding distributions.
\end{remark}

\section{Proofs}

\begin{proof}[Proof of Proposition~\ref{prop:exact}]
Let
\[
Z := \frac{S_a - \delta_D}{\sigma_D}
\quad\text{and}\quad
W_i := \frac{S_{v_i}}{\sigma_D},
\qquad i=1,\dots,M-1.
\]
Then
\[
Z\sim \mathcal N(0,1),
\qquad
W_1,\dots,W_{M-1}\stackrel{\mathrm{iid}}{\sim}\mathcal N(0,1),
\]
and $Z$ is independent of the $W_i$.

Retrieval succeeds exactly when
\[
S_a > \max_{1\le i\le M-1} S_{v_i},
\]
which is equivalent to
\[
Z + \Delta_D > \max_{1\le i\le M-1} W_i.
\]
Conditioning on $Z=x$ gives
\begin{align*}
P(\hat a=a)
&=
\int_{-\infty}^{\infty}
P\!\left(\max_{1\le i\le M-1} W_i < x+\Delta_D \,\middle|\, Z=x\right)\phi(x)\,dx \\
&=
\int_{-\infty}^{\infty}
P\!\left(\max_{1\le i\le M-1} W_i < x+\Delta_D\right)\phi(x)\,dx \\
&=
\int_{-\infty}^{\infty}
\prod_{i=1}^{M-1} P(W_i < x+\Delta_D)\,\phi(x)\,dx \\
&=
\int_{-\infty}^{\infty}
\Phi(x+\Delta_D)^{M-1}\phi(x)\,dx.
\end{align*}
This proves \eqref{eq:exact_integral}.
\end{proof}

\begin{lemma}[Gaussian maximum scale]\label{lem:max}
Let $U_M = \max\{W_1,\dots,W_{M-1}\}$ with $W_i\stackrel{\mathrm{iid}}{\sim}\mathcal N(0,1)$. If $M\to\infty$, then
\[
\frac{U_M}{\sqrt{2\log M}} \to 1
\qquad\text{in probability.}
\]
In particular, for every deterministic sequence $t_M$:
\begin{itemize}
    \item if $t_M - \sqrt{2\log M}\to +\infty$, then $P(U_M \le t_M)\to 1$;
    \item if $t_M - \sqrt{2\log M}\to -\infty$, then $P(U_M \le t_M)\to 0$.
\end{itemize}
\end{lemma}

\begin{proof}
This is a standard consequence of classical extreme-value theory for Gaussian maxima. One may use the more precise normalization
\[
U_M = a_M + b_M G_M,
\]
where
\[
a_M = \sqrt{2\log M} - \frac{\log\log M + \log(4\pi)}{2\sqrt{2\log M}},
\qquad
b_M = \frac{1}{\sqrt{2\log M}},
\]
and $G_M$ is tight. Since $a_M/\sqrt{2\log M}\to 1$ and $b_M\to 0$, the claimed convergence follows.
\end{proof}

\begin{proof}[Proof of Theorem~\ref{thm:threshold}]
Let
\[
U_D := \max_{1\le i\le M(D)-1} W_i.
\]
Then
\[
P(\hat a = a) = P(Z + \Delta_D > U_D),
\]
where $Z\sim\mathcal N(0,1)$ is independent of $U_D$.

We prove the two cases separately.

\medskip
\noindent\textbf{Case 1:} Assume $\Delta_D - L_D \to +\infty$, where $L_D=\sqrt{2\log M(D)}$.

Choose any deterministic sequence $r_D\to\infty$ such that
\[
r_D = o(\Delta_D - L_D).
\]
Then
\[
t_D := \Delta_D - r_D
\]
still satisfies
\[
t_D - L_D \to +\infty.
\]
By Lemma~\ref{lem:max},
\[
P(U_D \le t_D)\to 1.
\]
Also, since $r_D\to\infty$ and $Z\sim\mathcal N(0,1)$,
\[
P(Z > -r_D)\to 1.
\]
Now on the event $\{Z>-r_D\}\cap\{U_D\le t_D\}$ we have
\[
Z+\Delta_D > \Delta_D-r_D = t_D \ge U_D,
\]
so
\[
P(\hat a=a)
\ge
P(Z>-r_D,\; U_D\le t_D).
\]
Because $Z$ and $U_D$ are independent,
\[
P(\hat a=a)
\ge
P(Z>-r_D)\,P(U_D\le t_D)\to 1.
\]
Hence $P(\hat a=a)\to 1$.

\medskip
\noindent\textbf{Case 2:} Assume $\Delta_D - L_D \to -\infty$.

Choose any deterministic sequence $r_D\to\infty$ such that
\[
r_D = o(L_D-\Delta_D).
\]
Set
\[
s_D := \Delta_D + r_D.
\]
Then
\[
s_D - L_D \to -\infty.
\]
By Lemma~\ref{lem:max},
\[
P(U_D \le s_D)\to 0,
\qquad\text{equivalently}\qquad
P(U_D > s_D)\to 1.
\]
Also,
\[
P(Z > r_D)\to 0.
\]
Now if $Z+\Delta_D > U_D$, then either $Z>r_D$ or else $Z\le r_D$ and hence
\[
U_D < Z+\Delta_D \le \Delta_D + r_D = s_D.
\]
Therefore
\[
\{Z+\Delta_D > U_D\}
\subseteq
\{Z>r_D\}\cup\{U_D < s_D\},
\]
and so
\[
P(\hat a=a)
\le
P(Z>r_D) + P(U_D < s_D)\to 0.
\]
Thus $P(\hat a=a)\to 0$.
\end{proof}

\begin{proof}[Proof of Corollary~\ref{cor:hstar}]
Under $M \approx \bar d^{\,h}$,
\[
\log M \approx h\log\bar d,
\]
so the threshold scale from Theorem~\ref{thm:threshold},
\[
\Delta_D \approx \sqrt{2\log M},
\]
becomes
\[
\delta\sqrt{D} \approx \sqrt{2h\log\bar d}.
\]
Solving for $h$ yields
\[
h^*(D) \approx \frac{\delta^2 D}{2\log\bar d}.
\]
The one-sided asymptotic statements follow by substituting $\log M \approx h(D)\log\bar d$ into the hypotheses of Theorem~\ref{thm:threshold}.
\end{proof}

\section{Discussion}

Theorem~\ref{thm:threshold} isolates the basic mechanism in the Gaussian surrogate model: one signal score competes against the maximum of many noise scores, and the relevant scale for that maximum is $\sqrt{2\log M}$. The hop-count formula \eqref{eq:hstar} follows once the candidate-set size is approximated by $\bar d^{\,h}$.

The result has two important limitations.

First, it is a statement about the surrogate score model A2--A4. Real embedding systems may violate independence, equal variance, Gaussian tails, or the candidate-set growth law.

Second, the theorem does not imply that every retrieval algorithm obeys exactly the same threshold. Algorithms that alter the candidate set, use additional side information, exploit graph structure, or re-score candidates with a richer model may change the effective value of $M$, the score distribution, or the signal margin.

What the theorem does provide is a clean baseline: under the simplest iid Gaussian competition model, retrieval success is controlled by the gap between $\Delta_D$ and $\sqrt{2\log M(D)}$.

\end{document}